\scanpage{616}\mspage{4}
% archive sheet number: 615

avec les \deleted{\unsure{im}}\unsure{implications|inclusions}
tautologiques
\[
\begin{aligned}
\Phi\subset
\underset{\substack{\Vert\\\text{Q-saturé à}\\\text{g. simple}}}{\Phi'}
\subset{}&
\underset{\substack{\Vert\\\text{Q-saturé}\\\text{à g. transfini}}}{\Phi''}
\subset
\underset{\substack{\text{Q-saturé}\\\text{à g. total}}}{\Phi'''}
\subset
\underset{\substack{\text{Q-clôture)}\\\text{à gauche}}}{\Phi^{\sim}}
\end{aligned}\tag*{(1.7)}
\]
et dualement pour $\Psi$. Entre $\unsure{\Phi|\Psi}$ et $\Psi'$
il y a encore un $\Phi'_0$
\[
 \Phi\subset\Phi'_0\subset\Phi',\tag*{(1.8)}
\]
qui est le Q-\uline{saturé pour cochangement de base} (contenant
\inserted{aussi}, par définition, tous les iso). Pour construire
tous ces saturés, on procède ainsi :

1) $\Phi'_0=$ ens des flèches $i$\inserted{$:A\unsure{-|\to}\unsure{D|B}$}
s'insérant dans un carré cocartésien
\diagram{
\begin{tikzcd}[ampersand replacement=\&]
\deleted{\unsure{B}}A_0\arrow[r,"i_0"]\arrow[d] \& B_0\arrow[d]\\
 A\arrow[r,"i"] \& \unsure{D|B}
\end{tikzcd}
},
avec $i_0\in\Phi$ \deleted{\illegible{1}} ou $i_0$ une identité
(auquel cas on trouve pour $i$ les iso).

2) $\Phi'=$ ens des flèches composées multiples
\[
 i_n\,i_{n-1}\,\cdots\,i_1\deleted{\unsure{\xi|\}}}\,,
\]
avec $i_\alpha\in\Phi'_0$ \deleted{\unsure{$\xi$}}.

3) $\Phi''=$ ens des composés transfinis de flèches $\in\Phi'_0$.

\deleted{\unsure{4)}\unsure{$\Phi$}} On a envie de poser
\[
 \begin{aligned}
 \Phi'''=\Phi''\unsure{\natural|{}^{\natural}|\sharp}\overset{\mathrm{def}}{=}
 &\ \text{ens. des facteurs directs}\\
 &\ \text{d'éléments de }\Phi'',
 \end{aligned}
\]
et c'est un fait que $\Phi''\natural\subset\Phi'''$ trivialement, mais
il n'est pas clair en général qu'on ait
